% ============================================================================
% CHAPTER 9 — The 10-Step Process: From Blank Page to Standing Ovation
% Storymakers Method Report
% ============================================================================

\section{The 10-Step Process: From Blank Page to Standing Ovation}

\pullquote{Most presenters start by opening PowerPoint. That is like an architect starting by choosing doorknobs. The structure must exist before the surface does.}{Storymakers Method}

\sourcefig{infographics/ch09-10steps.jpg}{The 10-step process from blank page to standing ovation: Phase 0 (Launchpad), Phase 1 (Block level), Phase 2 (Loop level), and Phase 3 (Slide level).}

\subsection{The architecture of a presentation}

Every presentation that works---every pitch that closes, every keynote that resonates, every board deck that shifts opinion---follows the same underlying structure. Not the same story. Not the same slides. The same \emph{architecture}.

The Storymakers Method encodes that architecture into a 10-step process organised across four phases. The phases are not arbitrary. They map to the way human cognition processes persuasive communication: first the frame, then the argument, then the rhythm, then the signal.\endnote{Minto, Barbara, \emph{The Minto Pyramid Principle: Logic in Writing, Thinking, and Problem Solving}, Minto International, 1987. Minto's work at McKinsey established the foundational principle that structured communication begins with the answer, not the data.}

\begin{thesisbox}[The 10-Step Architecture]
\begin{description}[leftmargin=1.5em, labelindent=0em]
  \item[Phase 0 — Launchpad] (Step 0): SCQA, Big Idea, Pillars
  \item[Phase 1 — Block Level] (Steps 1--3): Structure, Journey, Evidence
  \item[Phase 2 — Loop Level] (Steps 4--6): Sequence, Tension, Storyboard
  \item[Phase 3 — Slide Level] (Steps 7--9): Insight, Concept, Signal
\end{description}
\end{thesisbox}

The numbering begins at zero deliberately. Phase 0 is the most important phase in the process. It is also the one that most people skip entirely.

\sourcefig{infographics/ch09-time-allocation.jpg}{The 60/25/15 split that most presenters invert: 60\% on foundation and blocks, 25\% on loops, 15\% on slides. Swap these ratios and you are decorating a house with no walls.}

\subsection{The time allocation that most people get wrong}

Before walking through each phase, a number that will feel counterintuitive: \textbf{60\% of total effort should go into Phase 0 and Phase 1.} The foundation and block-level structure. Not the slides. Not the visuals. Not the animations.

\begin{table}[H]
\centering
\small
\rowcolors{2}{altrow}{white}
\begin{tabularx}{\textwidth}{>{\bfseries\color{heading}}l X >{\color{signalred}\bfseries}c}
\toprule
\rowcolor{tableheader}
\textcolor{white}{\textbf{Phase}} & \textcolor{white}{\textbf{What happens here}} & \textcolor{white}{\textbf{Time}} \\
\midrule
Phase 0 + 1 & Foundation and block-level structure & 60\% \\
Phase 2 & Loop-level sequencing, tension, storyboarding & 25\% \\
Phase 3 & Slide-level refinement, visualisation, polish & 15\% \\
\bottomrule
\end{tabularx}
\caption{Time allocation across the Storymakers process. The 60/25/15 split is the single most violated guideline in corporate presentations.}
\label{tab:time-allocation}
\end{table}

Most professionals invert this ratio. They spend 15\% thinking about what to say and 60\% adjusting fonts, aligning boxes, and choosing stock photography. The result is a beautifully formatted argument that nobody follows, because the argument was never properly constructed.\endnote{Duarte, Nancy, \emph{Slide:ology: The Art and Science of Creating Great Presentations}, O'Reilly Media, 2008. Duarte's research at her firm found that executives spend an average of 36 hours per presentation, with 70\% of that time on visual formatting rather than message architecture.}

\begin{alertbox}[THE INVERSION TRAP]
If you are spending more time choosing fonts than choosing your argument structure, you are building a house by decorating the living room before pouring the foundation. No amount of visual polish rescues a presentation whose core argument is unclear.
\end{alertbox}

% ============================================================================
\subsection{Phase 0 --- The Launchpad (Step 0)}

Phase 0 is a single step, but it consumes 20--30\% of total project time. It answers three questions that every subsequent step depends on:

\begin{enumerate}
  \item \textbf{What is the situation, complication, question, and answer?} (SCQA)
  \item \textbf{What is the Big Idea?} (One sentence that captures the entire presentation)
  \item \textbf{What are the Pillars?} (The 2--4 MECE arguments that support the Big Idea)
\end{enumerate}

\subsubsection{SCQA: The strategic frame}

The SCQA framework---Situation, Complication, Question, Answer---is the argument engine of the Launchpad. See Chapter 4 for the full SCQA framework, test questions, and worked examples. In the 10-step process, the SCQA serves as the filter: every slide must connect back to this frame, or it does not belong.

\subsubsection{The Big Idea}

The Big Idea is a single sentence that captures the entire presentation. It has two components: it must state \emph{what is at stake} and it must contain \emph{your point of view}.\endnote{Duarte, Nancy, \emph{Resonate: Present Visual Stories That Transform Audiences}, Wiley, 2010. Duarte defines the Big Idea as ``your unique point of view that has something at stake for the audience'' and positions it as the gravitational centre of any effective presentation.}

A Big Idea is not a topic. ``Digital transformation in hospitality'' is a topic. ``Hotels that fail to build a unified guest record will lose 40\% of their repeat guests to competitors who do'' is a Big Idea. The difference is that the Big Idea demands a response: agreement, disagreement, or curiosity. A topic demands nothing.

\subsubsection{Pillars: The MECE argument}

Pillars are the 2--4 top-level arguments that hold the Big Idea up. They must be MECE: Mutually Exclusive (no overlap between pillars) and Collectively Exhaustive (together they cover the full argument). See Chapter 6 for the Architect's full MECE toolkit, including diagnostic examples and the MECE test at every pyramid level. The Pillar Test: if you remove one pillar and the argument does not collapse, it is decoration---merge it or cut it.

% ============================================================================
\subsection{Phase 1 --- Block Level (Steps 1--3)}

Phase 1 introduces the Creative Trio: three perspectives that collaborate to build the presentation's block-level structure. Each perspective asks a different question about the same material.\endnote{The Creative Trio framework draws on cross-functional design thinking principles documented in Brown, Tim, \emph{Change by Design: How Design Thinking Transforms Organizations and Inspires Innovation}, Harper Business, 2009.}

\subsubsection{Step 1: The Architect --- Structure the Problem}

The Architect asks: \emph{Is the argument logically complete?}

This step takes the SCQA and Pillars from Phase 0 and builds them into Blocks---the major structural sections of the presentation. A typical presentation has 3--5 Blocks. Each Block maps to a Pillar or a necessary structural element (introduction, conclusion, call to action).

The Architect checks for MECE compliance: do the Blocks overlap? Are there gaps? Is there a Block that exists because the presenter likes the topic, rather than because the argument requires it?

\begin{diagnosisbox}
\textbf{The Architect's Diagnostic.} Take your Block outline and show it to someone who knows nothing about the topic. Can they follow the argument from Block titles alone? If they cannot, the structure is broken. No amount of content will fix a structural problem.
\end{diagnosisbox}

\subsubsection{Step 2: The Storyteller --- Define the Journey}

The Storyteller asks: \emph{Does the audience feel something?}

This step maps the emotional arc of the presentation. Where does tension rise? Where does it resolve? Where does the audience feel the weight of the problem before seeing the solution?\endnote{Zak, Paul J., ``Why Your Brain Loves Good Storytelling,'' \emph{Harvard Business Review}, October 2014. Zak's research demonstrated that narrative tension triggers oxytocin release, which increases empathy and willingness to cooperate---directly relevant to persuasive presentations.}

The Storyteller is not decorating the Architect's structure with anecdotes. The Storyteller is ensuring that the argument has a \emph{shape}---a rise and fall of energy that keeps attention alive. A presentation with perfect logic but no emotional arc will be admired and forgotten. A presentation with emotional resonance and solid logic will be remembered and acted upon.

\subsubsection{Step 3: The Designer --- Plan the Evidence}

The Designer asks: \emph{What will they see?}

This step plans the evidence strategy for each Block: which data, which examples, which visuals will make the argument tangible. The Designer does not open PowerPoint. The Designer decides what type of evidence each argument needs: a number, a chart, a case study, a comparison, a metaphor, or a physical demonstration.\endnote{Reynolds, Garr, \emph{Presentation Zen: Simple Ideas on Presentation Design and Delivery}, New Riders, 2008. Reynolds' work established the principle that evidence selection---deciding \emph{what} to show---precedes and outranks visual design---deciding \emph{how} to show it.}

\begin{keybox}[THE CREATIVE TRIO IN PHASE 1]
\begin{itemize}[leftmargin=1.2em]
  \item \textbf{Architect}: ``The argument requires these 4 Blocks in this order.''
  \item \textbf{Storyteller}: ``The tension peaks in Block 2. Block 3 provides resolution.''
  \item \textbf{Designer}: ``Block 1 needs a data chart. Block 2 needs a case study. Block 4 needs a single number.''
\end{itemize}
The three perspectives are sequential within Phase 1 but iterative across phases. A discovery by the Designer often forces the Architect to restructure.
\end{keybox}

\subsection{The War Room: Where Ideas Become Structure}

The bridge between abstract strategy and tangible structure is storyboarding --- a technique borrowed from the film industry, where Walt Disney's team pioneered the practice of planning every scene visually before committing to animation.\endnote{The storyboarding technique was developed at Walt Disney Studios in the early 1930s. Disney artists pinned drawings to boards to visualise the flow of a story before animating, saving enormous time and cost.}

\begin{keybox}[THE WAR ROOM METHOD]
\begin{enumerate}
  \item \textbf{Post your Big Idea and SCQA} on red post-its at the top of a whiteboard. This is your north star.
  \item \textbf{Silent brainstorm:} Every team member writes one idea per post-it --- a data point, an argument, a question, a potential slide title. No discussion yet. This prevents groupthink.
  \item \textbf{Cluster and organise (MECE):} Together, group related post-its. Name each group with a header post-it. Ensure the groups are MECE.
  \item \textbf{Sequence the story:} Arrange the groups in a logical narrative sequence. Which chapter comes first? What builds on what?
  \item \textbf{Create the storyboard:} Draw larger boxes on the board for each slide. Write a headline (action title) and sketch the visual. This is your rolling pack.
\end{enumerate}
This method is powerful because it is fast, flexible, and visual. You can move slides, rewrite headlines, and restructure the story in real time.
\end{keybox}

\subsection{The Rolling Pack: Your Living Document}

The output of the storyboarding session is not a finished presentation. It is a \textbf{rolling pack} --- a low-fidelity PowerPoint or document that contains the complete structure but none of the polish. Every slide has a clear action title; in place of finished charts, there are descriptions: ``[CHART: Q1 vs Q2 sales comparison by region].''

\begin{alertbox}[THE FEEDBACK TRAP]
A team of junior consultants spent two weeks building an 80-slide presentation with complex analysis and polished charts. When the senior partner reviewed it, she said: ``The analysis is impressive, but the story is wrong. We're not answering the client's key question.'' The team had to start over. If they had shared a 10--15 slide rolling pack in the first three days, they would have received that feedback in time and saved two weeks of wasted work.\endnote{Composite case drawn from consulting practice. The rolling pack is standard methodology at McKinsey, BCG, and Bain for exactly this reason.}
\end{alertbox}

% ============================================================================
\subsection{Phase 2 --- Loop Level (Steps 4--6)}

If Blocks are chapters, Loops are paragraphs. A Loop is a self-contained argumentative unit within a Block: a single claim supported by evidence and followed by a transition. A typical Block contains 3--7 Loops.\endnote{The Loop concept as a unit of argumentation draws on Toulmin, Stephen E., \emph{The Uses of Argument}, Cambridge University Press, 1958. Toulmin's model of claim-data-warrant maps directly to the Loop's structure of assertion-evidence-transition.}

\subsubsection{Step 4: Sequence the Argument}

Each Loop within a Block must follow a logical sequence. The most common sequences are:

\begin{itemize}
  \item \textbf{Problem $\rightarrow$ Cause $\rightarrow$ Solution} (diagnostic narrative)
  \item \textbf{Claim $\rightarrow$ Evidence $\rightarrow$ Implication} (analytical narrative)
  \item \textbf{Before $\rightarrow$ After $\rightarrow$ How} (transformation narrative)
  \item \textbf{General $\rightarrow$ Specific $\rightarrow$ General} (inductive/deductive hybrid)
\end{itemize}

The choice of sequence is not aesthetic. It depends on the audience's starting position. An audience that already agrees with the problem needs a solution-first sequence. An audience that doubts the problem exists needs a diagnostic sequence.\endnote{Cialdini, Robert B., \emph{Influence: The Psychology of Persuasion}, revised edition, Harper Business, 2006. Cialdini's framework of pre-suasion---shaping the audience's mental frame before presenting the argument---directly informs Loop sequencing strategy.}

\subsubsection{Step 5: Manage Tension}

Tension is the engine of attention. A presentation without tension is a lecture. The Storyteller returns in this step to map the tension profile across Loops within each Block.

\begin{insightbox}[The Tension Principle]
Every Loop should contain a moment where the audience does not yet know the answer. This gap between question and resolution is what holds attention. If you reveal the answer before the audience feels the question, you have killed the tension. If you sustain the question too long, you have lost the audience. The window is narrow: 30 to 90 seconds of sustained uncertainty per Loop.
\end{insightbox}

\subsubsection{Step 6: Storyboard the Data}

Every chart, every number, every data visualisation in the presentation is planned in this step---not as a finished visual, but as a specification. The Designer returns to assign each Loop its evidence format.\endnote{Knaflic, Cole Nussbaumer, \emph{Storytelling with Data: A Data Visualization Guide for Business Professionals}, Wiley, 2015. Knaflic's framework for choosing chart types based on the \emph{relationship} you want to show (comparison, composition, distribution, relationship) is the standard reference for Step 6.}

The storyboard answers: What data supports this Loop? What chart type communicates the insight? What is the ``so what?'' of this chart? If you cannot state the ``so what'' in one sentence, the chart does not belong.

% ============================================================================
\subsection{Phase 3 --- Slide Level (Steps 7--9)}

Phase 3 is where PowerPoint (or Keynote, or Google Slides) finally opens. Not before. Only after 85\% of the intellectual work is complete.

\subsubsection{Step 7: Sharpen the Insight}

Every slide gets an \textbf{action title}---a complete sentence that states the slide's argument, not a topic label. See Chapter 6 for the Architect's full treatment of action titles versus topic titles, and Chapter 5 for the Headline Test. The core principle: if you read only the action titles from first slide to last, the entire argument should be clear without seeing any slide body.

\subsubsection{Step 8: Visualise the Concept}

Each slide's body must reinforce its action title with a single visual concept: a chart, a diagram, a photograph, a comparison, or a single number rendered large. The rule is one concept per slide.\endnote{Reynolds, Garr, \emph{Presentation Zen}, op.\ cit. Reynolds' ``one idea per slide'' principle draws on cognitive load theory: working memory can hold 4$\pm$1 chunks simultaneously (Cowan, Nelson, ``The Magical Number 4 in Short-Term Memory,'' \emph{Behavioral and Brain Sciences}, 2001).}

\begin{alertbox}[THE COGNITIVE LIMIT]
Working memory holds approximately four items simultaneously. A slide with six bullet points, a chart, and a footnote is asking the audience to process ten concepts at once. The result is not comprehension. It is noise. One concept per slide is not minimalism---it is cognitive science.
\end{alertbox}

\subsubsection{Step 9: Refine the Signal}

The final step removes everything that does not serve the argument. Every word, every colour, every line on a chart must justify its existence. This is where the Designer becomes an editor: cutting decoration, reducing labels, eliminating redundancy, and ensuring that the visual hierarchy on each slide matches the intellectual hierarchy of the argument.

Signal-to-noise ratio is the metric. A slide with 80\% signal and 20\% noise is good. A slide with 50\% signal and 50\% noise is mediocre. Most corporate slides are below 30\% signal.\endnote{Tufte, Edward R., \emph{The Visual Display of Quantitative Information}, 2nd edition, Graphics Press, 2001. Tufte's concept of ``data-ink ratio''---the proportion of a chart's ink devoted to displaying data versus non-data elements---is the quantitative foundation for the signal-to-noise principle.}

% ============================================================================
\subsection{Quality gates between phases}

Each phase boundary is a quality gate. You do not proceed to the next phase until the current phase passes its gate.

\begin{table}[H]
\centering
\small
\rowcolors{2}{altrow}{white}
\begin{tabularx}{\textwidth}{>{\bfseries\color{heading}}l X}
\toprule
\rowcolor{tableheader}
\textcolor{white}{\textbf{Gate}} & \textcolor{white}{\textbf{Pass criteria}} \\
\midrule
Phase 0 $\rightarrow$ 1 & SCQA written. Big Idea is one sentence with a point of view. Pillars are MECE and pass the ``removal test.'' \\
Phase 1 $\rightarrow$ 2 & Blocks defined. Emotional arc mapped. Evidence types assigned to every Block. \\
Phase 2 $\rightarrow$ 3 & Loops sequenced within every Block. Tension peaks identified. Data storyboarded with ``so what'' for every chart. \\
Phase 3 $\rightarrow$ Delivery & Every slide has an action title. One concept per slide. Title test passes (reading titles alone tells the story). \\
\bottomrule
\end{tabularx}
\caption{Quality gates between phases. Skipping a gate is the most common source of presentation failure.}
\label{tab:quality-gates}
\end{table}

\begin{insightbox}[Why Gates Matter]
The cost of fixing a structural problem increases exponentially as you move through phases. A Pillar change in Phase 0 costs 30 minutes. The same change discovered in Phase 3---after slides are built---costs 3 hours. Quality gates exist to catch structural flaws early, when they are cheap to fix.
\end{insightbox}

% ============================================================================
\subsection{The process is iterative, not linear}

The 10 steps are numbered for clarity, not because they execute in strict sequence. Every practitioner of the method will revisit earlier steps as new insights emerge in later steps. This is not failure. It is the process working correctly.

Common iteration patterns:

\begin{itemize}
  \item \textbf{Step 6 $\rightarrow$ Step 1}: Storyboarding data reveals that a Block lacks sufficient evidence. The Architect restructures.
  \item \textbf{Step 7 $\rightarrow$ Step 2}: Writing action titles reveals that the emotional arc is flat. The Storyteller remaps.
  \item \textbf{Step 8 $\rightarrow$ Step 0}: Visualising the concept exposes a weakness in the Big Idea itself. Back to the Launchpad.
  \item \textbf{Step 5 $\rightarrow$ Step 4}: Managing tension within a Loop reveals that the Loop sequence is wrong. Reorder.
\end{itemize}

\begin{diagnosisbox}
\textbf{The Iteration Signal.} If you never revisit an earlier step, you are either a genius or you are not looking hard enough. In practice, teams iterate 2--3 times across phases before converging. The process is a spiral, not a waterfall. Each iteration tightens the argument, sharpens the story, and clarifies the evidence.
\end{diagnosisbox}

\pullquote{A presentation is not built in ten steps. It is built in ten steps, revisited three times, with the courage to kill ideas that felt brilliant at Step 2 but collapse at Step 7.}{Storymakers Method}

\clearpage
